\documentclass[12pt,a4paper]{article}

% Codificación y fuentes
\usepackage[utf8]{inputenc}
\usepackage[T1]{fontenc}
\usepackage{lmodern}

% Idioma español
\usepackage[spanish,es-noquoting,es-tabla]{babel}
% es-noquoting: para usar comillas inglesas " " sin conflicto
% es-tabla: cambia "Cuadro" por "Tabla"

% Matemáticas y símbolos
\usepackage{amsmath,amssymb,amsfonts}
\usepackage{amsthm}

% Gráficos y colores
\usepackage{graphicx}
\usepackage{xcolor}
\usepackage{booktabs} % Tablas profesionales

% Geometría de página
\usepackage[margin=2.5cm]{geometry}

% Hiperenlaces
\usepackage[colorlinks=true,linkcolor=blue,citecolor=blue,urlcolor=blue]{hyperref}

% Configuración de párrafos
\usepackage{setspace}
\onehalfspacing % Interlineado 1.5
\setlength{\parskip}{0.5em}

% Teoremas y definiciones (opcional)
\theoremstyle{definition}
\newtheorem{definicion}{Definición}[section]
\theoremstyle{theorem}
\newtheorem{teorema}{Teorema}[section]
\newtheorem{proposicion}{Proposición}[section]

% Metadatos
\title{\textbf{Título del Artículo}}
\author{Nombre del Autor \\
	\small Universidad/Institución \\
	\small \texttt{correo@ejemplo.com}}
\date{\today}

\begin{document}

\maketitle

\begin{abstract}
\noindent Escribe aquí el resumen de tu artículo. Debe ser conciso, informativo y reflejar el contenido principal del trabajo. Generalmente entre 150-250 palabras.

\vspace{1em}
\noindent\textbf{Palabras clave:} keyword1, keyword2, keyword3, keyword4.
\end{abstract}

\section{Introducción}

Comienza con el contexto general y la motivación de tu trabajo. Puedes citar referencias así: \cite{autor2023} o \cite{autor2023} dependiendo del estilo.

\subsection{Antecedentes}

Desarrolla el marco teórico necesario.

\subsubsection{Subsección de ejemplo}

Texto de la subsección.

\section{Metodología}

Describe los métodos utilizados. Por ejemplo, una ecuación matemática:
\begin{equation}
	E = mc^2
	\label{eq:einstein}
\end{equation}

En la ecuación \eqref{eq:einstein} se muestra la equivalencia masa-energía.

\section{Resultados}

\begin{table}[htbp]
	\centering
	\caption{Ejemplo de tabla profesional}
	\begin{tabular}{@{}lcc@{}}
		\toprule
		Concepto & Valor A & Valor B \\
		\midrule
		Elemento 1 & 10.5 & 20.3 \\
		Elemento 2 & 15.2 & 18.7 \\
		\bottomrule
	\end{tabular}
	\label{tab:ejemplo}
\end{table}

\section{Conclusiones}

Resume los hallazgos principales y su importancia.

\section*{Agradecimientos}

Opcional: agradecimientos a financiaciones o colaboradores.

% Bibliografía
\begin{thebibliography}{99}

\bibitem{autor2023}
Apellido, N. (2023).
\newblock Título del libro o artículo.
\newblock \textit{Revista}, 12(3), 45--67.

\bibitem{otro2022}
Otro Autor, M. (2022).
\newblock \textit{Título en cursiva}.
\newblock Editorial.

\end{thebibliography}

% O usa BibTeX/BibLaTeX:
% \bibliographystyle{apalike-es}
% \bibliography{referencias}

\end{document}
